% Part: first-order-logic
% Chapter: syntax-and-semantics
% Section: covered-structures

\documentclass[../../../include/open-logic-section]{subfiles}

\begin{document}

\olfileid[pt]{fol}{syn}{cov}

\olsection{\printtoken{P}{structure} Cobertas para Linguagens de Primeira Ordem}

\begin{explain}
Lembre-se de que um termo é \emph{fechado} se não contém !!{variable}s.
\end{explain}

\begin{defn}[!!^{value} de termos fechados]
Se $t$ é um termo fechado da linguagem~$\Lang L$ e $\Struct M$ é
!!a{structure} para~$\Lang L$, o \emph{!!{value}}~$\Value{t}{M}$ é
definido como segue:
\begin{enumerate}
\item Se $t$ é apenas o !!{constant}~$c$, então $\Value{c}{M} = \Assign{c}{M}$.
\item Se $t$ é da forma $\Atom{f}{t_1, \ldots, t_n}$, então
  \[
  \Value{t}{M} = \Assign{f}{M}(\Value{t_1}{M}, \ldots,
  \Value{t_n}{M}).
  \]
\end{enumerate}
\end{defn}

\begin{defn}[!!^{structure} Coberta]
!!^a{structure} é \emph{coberta} se todo elemento do domínio é o
!!{value} de algum termo fechado.
\end{defn}

\begin{ex}
Seja~$\Lang L$ a linguagem com os !!{constant}s $\Obj{zero}$,
$\Obj{one}$, $\Obj{two}$, \dots, o !!{predicate} binário~$<$, e os
!!{function}s binários $+$ e $\times$. Então !!a{structure}~$\Struct
M$ para~$\Lang L$ é aquela com domínio $\Domain M = \{0, 1, 2, \ldots
\}$ e atribuições $\Assign{\Obj{zero}}{M} = 0$,
$\Assign{\Obj{one}}{M} = 1$, $\Assign{\Obj{two}}{M} = 2$, e assim
por diante. Para o símbolo de relação binário $<$, o conjunto $\Assign{<}{M}$ é
o conjunto de todos os pares $\tuple{c_1, c_2} \in \Domain{M}^2$ tais que
$c_1$ é menor que~$c_2$: por exemplo, $\tuple{1, 3} \in
\Assign{<}{M}$ mas $\tuple{2, 2} \notin \Assign{<}{M}$. Para o
!!{function} binário $+$, defina $\Assign{+}{M}$ da maneira usual---por
exemplo, $\Assign{+}{M}(2,3)$ resulta em~$5$, e analogamente para o
!!{function} binário~$\times$. Portanto, o !!{value} de $\Obj{four}$ é
apenas~$4$, e o !!{value} de $\times(\Obj{two},
+(\Obj{three},\Obj{zero}))$ (ou em notação infixa, $\Obj{two} \times
(\Obj{three} + \Obj{zero})$) é
\begin{multline*}
\Value{\times(\Obj{two}, +(\Obj{three},\Obj{zero}))}{M} =\\
\begin{aligned}
& =\Assign{\times}{M}(\Value{\Obj{two}}{M}, \Value{+(\Obj{three}, \Obj{zero})}{M})\\
& = \Assign{\times}{M}(\Value{\Obj{two}}{M}, \Assign{+}{M}(\Value{\Obj{three}}{M},
\Value{\Obj{zero}}{M})) \\
& = \Assign{\times}{M}(\Assign{\Obj{two}}{M}, \Assign{+}{M}(\Assign{\Obj{three}}{M},
\Assign{\Obj{zero}}{M})) \\
& = \Assign{\times}{M}(2, \Assign{+}{M}(3, 0)) \\
& = \Assign{\times}{M}(2, 3) \\
& = 6
\end{aligned}
\end{multline*}
\end{ex}

\begin{prob}
$\Struct N$, o modelo padrão da aritmética, é coberto? Explique.
\end{prob}

\end{document}
